%%%%%%%%%%%%%%%%%%%%%%%%%%%%%%%%%%%%%%%%%%%%%%%%%%%%%%%%%%%%%%%%%%%%%%
% Lớp 11 | Quan hệ vuông góc trong không gian | Hai mặt phẳng vuông góc
% Dạng 1H8?4-4: Dựng mặt phẳng, thiết diện
% Số câu: 9  (sinh tự động bởi tach_cau.py; từ nay thầy sửa trực tiếp file này)
%%%%%%%%%%%%%%%%%%%%%%%%%%%%%%%%%%%%%%%%%%%%%%%%%%%%%%%%%%%%%%%%%%%%%%

%%%%%%%%%%%%%%%% TRẮC NGHIỆM 4 PHƯƠNG ÁN %%%%%%%%%%%%%%%%

% nguồn: Ok_Toan11_TracNghiem_Lan2
\begin{ex}%[1H8V4-4]%[ma:54690d8f]
	 Cho hình chóp $S.ABCD$ có đáy $ABCD$ là hình chữ nhật tâm $O$ với $AB=a$; $AD=2a$. Cạnh bên $SA=a$ và vuông góc với đáy. Gọi ($\alpha$) là mặt phẳng qua $SO$ và vuông góc với $(SAD)$. Tính diện tích $S$ của thiết diện tạo bởi $(\alpha)$ và hình chóp đã cho
	 \choice
		 {$S=\dfrac{a^2\sqrt{3}}{2}$}
		 {\True$S=\dfrac{a^2\sqrt{2}}{2}$}
		{$S=\dfrac{a^2}{2}$} 
		{$a^2$} 
	\loigiai{ 
		\begin{center}
			\begin{tikzpicture}[scale=1,font=\footnotesize,line join=round,line cap=round,>=stealth]
					\path 
						(0,0) coordinate (A) 
						 ++(-140:2.5) coordinate (B)
						 ++(0:5) coordinate (C)
						($(A)+(C)-(B)$)  coordinate (D)
						($(A)+(0,3)$) coordinate (S) 
						($(A)!0.5!(D)$) coordinate (M) 
						($(B)!0.5!(C)$) coordinate (N) 
						 ($(A)!0.5!(C)$) coordinate (O);
					\fill[blue!30]  (S)--(M)--(N);
					\foreach \i in{B,C,D}{\draw (S)--(\i);}; 
					\draw  (B)--(C)--(D)
						 (S)--(N);
					\draw[dashed] (S)--(A)--(B) 
						 (A)--(D)
						(A)--(C) 
						(B)--(D) 
						(S)--(M)--(N) 
						(S)--(O); 
					\pic[draw,angle  eccentricity=1.8,angle radius=2mm]{right angle=S--A--D};
					\foreach \i/\g  in {A/180,B/-90,C/-90,D/0,S/90,M/30,N/-90,O/-65}
					\fill (\i) circle  (1pt)+(\g:3mm) node[scale=1]{$\i$};
				\end{tikzpicture}
		\end{center}
		 Gọi $M$; $N$ lần lượt là trung điểm $AD$; $BC$.\\
		Khi đó $MN$ đi qua $O$ và $\left\{\begin{array}{l}MN\perp AD\\ MN\perp SA
	\end{array} \Rightarrow MN\perp(SAD)\right.$\\
	Từ đó suy ra $(\alpha) \equiv(SMN)$ và thiết diện cần tìm là tam giác  $SMN$.\\
	Tam  giác $SMN$ vuông tại $M$
	 nên $S_{\triangle SMN}=\dfrac{1}{2} SM\cdot MN=\dfrac{1}{2} \sqrt{SA^2+\left(\dfrac{AD}{2}\right)^2} \cdot AB=\dfrac{a^2\sqrt{2}}{2}$.
	 }
\end{ex}

% nguồn: Ok_Toan11_TracNghiem_Lan2
\begin{ex}%[1H8V4-4]%[ma:7b8fe05c]
	Cho hình chóp tam giác đều $S. ABC$ đỉnh $S$, có độ dài cạnh đáy bằng $a$. Gọi $M$ và $N$ lần lượt là trung điểm của các cạnh $SB$ và $SC$. Biết mặt phẳng $(AMN)$ vuông góc  với mặt phẳng $(SBC.$. Tính diện tích tam giác $AMN$ theo $a$.
	 \choice
		 {$\dfrac{a^2\sqrt{10}}{24}$}
		{$\dfrac{a^2\sqrt{10}}{16}$} 
		 {$\dfrac{a^2\sqrt{5}}{8}$}
		 {$\dfrac{a^2\sqrt{5}}{4}$}
	 \loigiai{
		\begin{center}
			\begin{tikzpicture}[scale=1,font=\footnotesize,line join=round,line cap=round,>=stealth]
					 \path
						 (0,0) coordinate (A)
						++(-30:3) coordinate (B) 
						(4,0) coordinate  (C)
						($(B)!1/2!(C)$)  coordinate (I)
						($(A)!2/3!(I)$) coordinate (O) 
						($(O)+(0,3)$) coordinate  (S)
						($(S)!1/2!(B)$) coordinate  (M)
						($(S)!1/2!(C)$) coordinate  (N)
						($(M)!1/2!(N)$) coordinate  (H);
					\foreach \i in{A,B,C}{\draw (S)--(\i);}; 
					\draw (A)--(B)--(C) 
						 (A)--(M)--(N)
						 (S)--(I);
					\draw[dashed]  (A)--(C)
						 (A)--(I)
						(S)--(O) 
						 (A)--(H)
						 (A)--(N);
					 \pic[draw,angle eccentricity=1.8,angle radius=2mm]{right angle=S--O--A};
					\foreach \i/\g in  {S/90,A/180,C/0,B/-90,I/-90,O/-90,M/-40,N/20,H/-20}
					 \fill (\i) circle (1pt)+(\g:4mm) node[scale=1]{$\i$};
				\end{tikzpicture}
		\end{center}
		Ta thấy do hình chóp $S.ABC$ đỉnh  $S$ là chóp tam giác đều nên $AB=BC=AC=a$.
		$$ 
		\triangle  S A B=\triangle S A C(\text{c.c.c}) \Rightarrow A M=A N.
		$$ 
		Do đó tam  giác $AMN$ cân tại $A$. Gọi $H$ là trung điểm của $MN$ thì $AH\perp MN$ và $I$ là trung điểm của $BC$.\\
		 $\heva{&(AMN) \perp(SBC) \\
		&(AMN) \cap(SBC)=MN\\&AH\perp MN } \Rightarrow A H \perp S H; A H \perp S  I $.\\
		Xét tam giác $SAI$  có $AH$ vừa là trung tuyến vừa là đường cao nên tam giác $SAI$ cân tại $A$.\\
		Tam giác $ABC$  đều cạnh $a \Rightarrow AI=\dfrac{a \sqrt{3}}{2}=SA=SB$.\\
		Ta có: $S H=\dfrac{1}{2} S  I=\dfrac{a}{2 \sqrt{2}}$.\\
		Xét tam giác $A S H$  vuông tại $H$ nên $A H=\sqrt{S A^2-S H^2}=\sqrt{\dfrac{3 a^2}{4}-\dfrac{a^2}{8}}=\dfrac{a \sqrt{5}}{2 \sqrt{2}}$.\\
		Vậy $S_{A  M N}=\dfrac{1}{2} \cdot A H \cdot M N=\frac{1}{2} \cdot \dfrac{a \sqrt{5}}{2 \sqrt{2}} \cdot \dfrac{a}{2}=\dfrac{a^2 \sqrt{10}}{16}$.
	 }
\end{ex}


%%%%%%%%%%%%%%%% ĐÚNG/SAI %%%%%%%%%%%%%%%%

% nguồn: Ok_dungsai_lop11_Lan1
\begin{ex}%[1H8V4-4]%[ma:418270d5]
  Cho hình chóp $S.ABCD$ có đáy $ABCD$ là hình thang vuông tại $A$, đáy lớn $AD=8$, đáy bé $BC=6$; cạnh $SA$ vuông góc với mặt phẳng $(ABCD)$ và $SA=6$. Gọi $M$ là trung điểm cạnh $AB$ và $(P)$ là mặt phẳng qua $M$, vuông góc với $AB$. Xét tính Đúng - Sai của các mệnh đề sau:
 \choiceTF{$CD \perp  (P)$}
  {\True $(P)  \parallel (SAD)$}
  {\True Thiết diện của $(P)$ và hình chóp  là hình thang vuông}
  {Diện tích thiết diện  của $(P)$ và hình chóp là $30$}
 \loigiai{ 
  \begin{center}
   \begin{tikzpicture}[scale=0.7]
     \coordinate  (a) at (0,0);
     \coordinate (b)  at (-2,-3);
      \coordinate (c) at (2,-3);
     \coordinate (d) at (6,0); 
      \coordinate (s) at ($(a)+(0,4)$);
      \coordinate (m) at ($(a)!.5!(b)$);
     \coordinate (n)  at ($(c)!.5!(d)$);
     \coordinate (p) at  ($(s)!.5!(c)$);
      \coordinate (q) at ($(s)!.5!(b)$);
     \draw  [dashed] (a) node[left]{$A$}--(b) node[below]{$B$}
      (a)--(d) node[right]{$D$} 
      (s) node[above]{$S$}--(a); 
      \draw (b)--(c) node[below]{$C$}
       (s)--(b)
      (s)--(c)--(d)--cycle (q)--(p) node[above  right]{$P$}--(n);
     \draw [dashed] (q) node[left]{$Q$}--(m) node[below right]{$M$}--(n)  node[right]{$N$};
     \draw pic[draw=black,angle radius=0.2cm] {right angle =  b--a--d};
     \draw pic[draw=black,angle radius=0.2cm] {right angle = a--b--c}; 
     \draw pic[draw=black,angle radius=0.2cm]  {right angle = q--m--n};
     \foreach \x in{s,a,b,c,d,m,n,p,q}{ 
       \fill (\x) circle (1pt);
      }
    \end{tikzpicture}
  \end{center}
  \begin{itemchoice}
   \itemch  Sai.\\
    Giả sử $CD \perp (P)$, mà ta có $AB \perp (P)$ nên suy ra $AB \parallel CD$. Mà $AB \nparallel CD$. Do đó đây là khẳng định sai.
   \itemch  Đúng.\\
   Ta  có $\heva{&AB \perp (SAD) \\& AB \perp (P)} \Rightarrow (P) \parallel (SAD)$.
   \itemch  Đúng.\\
   - Từ $\heva{&AB\perp AD\\ &AB\perp SA}\Rightarrow  AB\perp (SAD)$, suy ra $(P)\parallel (SAD)$.\\
   - Từ  đó dựng được thiết diện là hình thang $MNPQ$.\\
   - Dễ thấy $MNPQ$ là hình thang vuông tại $M$  và $Q$.
    \itemch Sai.\\
    - Tính được: $MN=7$, $MQ=3=PQ$.\\
   - Diện  tích thiết diện là
    $$S=\dfrac{(7+3)\cdot 3}{2}=15.$$
  \end{itemchoice}
 } 
\end{ex}

% nguồn: Ok_dungsai_lop11_Lan1
\begin{ex}%[1H8V4-4]%[ma:7fd2d5f0]
 Cho tứ diện đều $ABCD$ có cạnh bằng $a$. Gọi $M$, $N$ lần lượt là trọng tâm  của tam giác $ABC$ và tam giác $ ABD$.
  \choiceTF
  {\True Đường thẳng $AB$ và $CD$  chéo nhau}
  {\True  $MN$ song song với $CD$}
  {Thiết diện của hình tứ diện cắt bởi  $(BMN)$ là hình tứ giác}
  {\True Diện tích thiết diện của hình tứ diện cắt bởi  mặt phẳng $(BMN)$ bằng $\dfrac{a^2\sqrt{11}}{16}$}
 \loigiai{ 
  \begin{center}
   \begin{tikzpicture}[scale=1,font=\footnotesize,line join=round,line cap=round,>=stealth]
     \def\ac{4.5} 
     \def\ab{2} 
      \def\h{4}
     \def\gocA{45} 
     \coordinate[label=left:$A$] (A)  at (0,0);
     \coordinate[label=below left:$B$] (B) at  (-\gocA:\ab);
     \coordinate[label=right:$C$]  (C) at (\ac,0);
     \coordinate[label=below right:$E$] (E)  at ($(A)!.5!(C)$);
     \coordinate[label=below right:$M$] (M)  at ($(B)!2/3!(E)$);
     \coordinate[label=above left:$D$]  (D) at ($(M)+(90:\h)$);
      \coordinate[label=left:$F$] (F) at ($(D)!.5!(A)$);
     \coordinate[label=left:$N$] (N) at  ($(B)!2/3!(F)$);
     \draw  (A)--(B)--(C)--(D)--cycle (D)--(B)
       (B)--(F);
     \draw[dashed] (F)--(E)--(B) 
       (A)--(C)
      (N)--(M)--(D); 
      \foreach \diem in {A,B,C,E,M,D,F,N} \fill (\diem) circle (1pt);
    \end{tikzpicture}
  \end{center}
  Gọi $E$, $F$ lần lượt là trung điểm của các cạnh  $AC$, $AD$. Ta có
  \begin{itemchoice}
   \itemch $AB$ và $CD$ chéo nhau. 
   \itemch Xét tam giác $BEF$, có  $\dfrac{BN}{BF}=\dfrac{2}{3}=\dfrac{BM}{BE}$ nên $MN \parallel FE$.
    \itemch Thiết diện của hình tứ diện và $(BMN)$ là tam giác $BEF$.
   \itemch Diện tích thiết diện là diện tích tam giác  $BEF$. \\
   Xét tam giác $BEF$ có $BF=BE=\dfrac{a\sqrt{3}}{2}$, $MN=\dfrac{DC}{2}=\dfrac{a}{2}$.\\ 
   Ta có $\cos \widehat{EBF}=\dfrac{BF^2+BE^2-EF^2}{2\cdot  BF\cdot BE}=\dfrac{5}{6}$.\\
   Suy ra $\sin \widehat{EBF}=\dfrac{\sqrt{11}}{6}$.  Diện tích tam giác $BEF$ là
   $$S_{BEF}=\dfrac{1}{2}BE\cdot BF\cdot \sin \widehat{EBF}=\dfrac{\sqrt{11}a^2}{16}.$$ 
  \end{itemchoice}
 } 
\end{ex}

% nguồn: Ok_dungsai_lop11_Lan1
\begin{ex}%[1H8V4-4]%[1H8H2-2]%[ma:e79cd98c]
 Cho hình lăng trụ $A B C A' B' C'$ có đáy là tam giác giác vuông cân tại $A$ với cạnh huyền $B C=2 a$. Biết $A' H \perp(A B C)$ với $H$ là trung điểm  $B C$. Khi đó:
 \choiceTF 
  {\True $B C \perp\left(A A' H\right)$} 
  {\True $B' C' \perp  A A'$}
  {Tìm được hình chiếu của tam giác $A' A  B$ trên mặt phẳng $(A B C)$ khi đó, diện tích hình chiếu đó theo $a$ bằng: $\dfrac{a^2}{3}$}
  {Gọi $I$ là hình chiếu của $A'$ trên mặt phẳng $\left(B C C' B'\right)$. Biết $A' I=\dfrac{a}{2}$.  Khi đó độ dài $A' H$ theo $a$ bằng: $\dfrac{a \sqrt{3}}{2}$}
  \loigiai{
  \begin{center}
   \begin{tikzpicture}
      \def\a{3}
     \def\h{4.5} 
     \path (0:0) coordinate  (A)
       ++(0:\a) coordinate (C)
      ++(-135:\a/2) coordinate (B) 
      ($(B)!0.5!(C)$) coordinate (H) 
       ($(H)+(90:\h)$) coordinate (A')
      ($(A')+(C)-(A)$) coordinate  (C')
      ($(C')+(B)-(C)$) coordinate  (B')
      ($(B')!0.5!(C')$) coordinate  (K)
      ($(K)!1/4!(H)$) coordinate  (I);
     \draw[dashed,thick]  (A')--(H)--(A)--(C)
      (I)--(A') 
       (K)--(H);
     \draw[thick]  (C)--(C')
      (B)--(B') 
       (A)--(A')--(K)
      (A)--(B)--(C) 
       (A')--(B')--(C')--cycle;
     \foreach  \x/\g in {A/180,B/-45,C/0,A'/180,B'/90,C'/0,H/-45,K/0,I/0}
      \fill[black] (\x) circle (1pt)
       ($(\g:4mm)+(\x)$) node {$\x$};
      \draw pic[draw,angle radius=2mm]{right angle=A'--H--C};
      \draw pic[draw,angle radius=2mm]{right angle=A'--K--C'};
     \draw pic[draw,angle  radius=2mm]{right angle=A'--I--H};
    \end{tikzpicture}
  \end{center}
  \begin{itemchoice}
   \itemch Đúng.$\heva{&B C \perp A' H \,\,( \text{do }A' H \perp(A B C)) \\ &B C \perp A H \,\,( \text{do }\triangle A B C \text { vuông cân tại } A, H \text { là trung điểm của }  B C)}$ \\ $\Rightarrow B C \perp(A A' H)$
   \itemch Đúng. Do $B'  C' \parallel B C$ nên $B' C' \perp\left(A A' H\right) \Rightarrow B' C' \perp A A'$.
   \itemch Sai. Vì $A' H \perp(A B C)$ nên  hình chiếu của $A A'$ trên $(A B C)$ là $A H$, hình chiếu của $A' B$ trên $(A B C)$ là $B H$.\\
   Vậy hình chiếu của tam giác $A' A B$  trên mặt phẳng $(A B C)$ chính là tam giác $A B H$.\\
   Tam giác $A B C$ vuông cân  tại $A$ có $B C=2 a \Rightarrow A B=A C=a \sqrt{2}$.\\
   Diện tích tam  giác $A B H$ là:
   $S_{\triangle A B H}=\dfrac{1}{2} S_{A B C}=\dfrac{1}{2} \cdot \dfrac{1}{2}  A B \cdot A C=\dfrac{1}{4}(a \sqrt{2})^2=\dfrac{a^2}{2}$
   \itemch Sai. Gọi $K$ là trung điểm $B' C'$. Dễ thấy $\left(A' A H\right) \equiv\left(A H K  A'\right)$.
   Mà $B' C' \perp\left(A  A' H\right)$ nên $B' C' \perp\left(A H K A'\right)$.
   Trong  mặt phẳng $\left(A H K A'\right)$, kẻ $A' I \perp H K$ tại $I$. (1)\\
   Vì $B' C' \perp\left(A H K A'\right), A' I \subset\left(A H K  A'\right)$ nên $A' I \perp B' C'$. (2)\\
   Từ $(1)$ và $(2)$ suy ra $A' I \perp\left(B C C' B'\right)$ hay $I$ là hình  chiếu của $A'$ trên mặt phẳng $\left(B C C' B'\right)$.\\
   Tam giác $A' B' C'$ vuông cân tại  $A'$ nên $A' K=\dfrac{B' C'}{2}=a$.\\
   Tam giác $A' H K$ vuông tại $A'$ có đường cao  $A' I$ nên ta có:\\
   $\dfrac{1}{A' I^2}=\dfrac{1}{A' H^2}+\dfrac{1}{A' K^2} \Rightarrow \dfrac{1}{A' H^2}=\dfrac{1}{\dfrac{a^2}{4}}-\dfrac{1}{a^2}=\dfrac{3}{a^2} \Rightarrow A' H=\dfrac{a \sqrt{3}}{3}$ 
  \end{itemchoice}
 } 
\end{ex}


%%%%%%%%%%%%%%%% TỰ LUẬN %%%%%%%%%%%%%%%%

% nguồn: Ok_TuanLuan_Toan11_Lan2
\begin{ex}%[1H8H4-4]%[ma:3fa09265]
	Cho hình chóp $S.ABCD$ có đáy là hình thang vuông tại $A,  D$; $AB=2a,\ SA=AD=DC=a$; $SA\perp (ABCD)$. Tính diện tích thiết diện tạo bởi mặt phẳng $(\alpha)$ qua $SD$ và $(\alpha)\perp (SAC)$.
	 \loigiai{
		\begin{center}
			\begin{tikzpicture}[line cap=round,line join=round,>=triangle 45,x=1cm,y=1cm]
					 \draw (0,4)-- (0,0);
					\draw (0,0)--  (1,-1);
					\draw (1,-1)--  (4,-1);
					\draw  (4,-1)-- (6,0);
					\draw  [dashed] (6,0)-- (0,0);
					\draw (1,-1)--  (0,4);
					\draw (0,4)-- (4,-1); 
					\draw  (6,0)-- (0,4);
					 \draw [dashed] (4,-1)-- (3,0);
					\draw  [dashed] (3,0)-- (1,-1);
					\draw [dashed] (0,0)-- (4,-1); 
					 \draw [dashed] (2,-0.5)-- (0,4);
					\draw [dashed] (3,0)--  (0,4);
					\draw  [fill] (0,0) circle (1pt);
					\draw[color=black] (-0.28,0.14) node  {$A$};
					\draw [fill] (6,0) circle (1pt); 
					 \draw[color=black] (6.3,0.38) node {$B$};
					\draw [fill] (4,-1) circle (1pt); 
					\draw[color=black]  (4.06,-1.2) node {$C$};
					\draw  [fill] (1,-1) circle (1pt);
					\draw[color=black] (0.7,-1.12)  node {$D$};
					\draw [fill] (0,4) circle (1pt); 
					 \draw[color=black] (-0.18,4.54) node {$S$};
					\draw [fill]  (3,0) circle (1pt);
					\draw[color=black] (2.42,0.36) node {$M$}; 
					\draw [fill]  (2,-0.5) circle (1pt);
					\draw[color=black] (2.1,-0.08) node {$O$}; 
				\end{tikzpicture}
		\end{center}
		Gọi  $M$ là trung điểm $AB$. \\
		 Tứ giác $ADCM$ là hình vuông $\Rightarrow DM\perp AC$. \\
		Mà $DM\perp  SA$ suy ra \\
		$DM\perp (SAC)\Rightarrow (SDM)\perp  (SAC)\Rightarrow (\alpha)=(SDM)$. \\
		Suy ra thiết diện là $\triangle SDM$.  \\
		Ta có $SO=\sqrt{SA^2+OA^2}=\dfrac{a\sqrt{6}}{2},\ DM=a\sqrt{2}$.  \\
		Diện tích thiết  diện là $S_{\triangle SDM}=\dfrac{SO\cdot DM}{2}=\dfrac{a^2\sqrt{3}}{2}$.
	} 
\end{ex}

% nguồn: Ok_TuanLuan_Toan11_Lan2
\begin{ex}%[1H8H4-4]%[ma:00eb55a9]
	Cho hình chóp $S.ABCD$ có đáy là hình vuông cạnh $a$, cạnh bên $SA$ vuông góc  với mặt phẳng đáy và $SA=2a$. Mặt phẳng $(P)$ qua $A$ và vuông góc với $SC$. Tính diện tích của thiết diện cắt bởi $(P)$ và hình chóp $S.ABCD$.
	\loigiai{ 
		\begin{center}
			\begin{tikzpicture}[line cap=round,line join=round,>=triangle 45,x=1cm,y=1cm]
					\draw [dashed] (0,0)--  (-1,-1);
					\draw  (-1,-1)-- (3,-1);
					 \draw (3,-1)-- (4,0);
					\draw [dashed]  (4,0)-- (0,0);
					 \draw (3,-1)-- (0,4);
					\draw (0,4)-- (-1,-1); 
					\draw  (4,0)-- (0,4);
					 \draw [dashed] (0,4)-- (0,0);
					\draw  [dashed] (0,0)-- (-0.5,1.5);
					\draw [dashed]  (-0.5,1.5)-- (2,2);
					\draw  [dashed] (2,2)-- (0,0);
					\draw (-0.5,1.5)-- (1,2.3333333333333335); 
					\draw (1,2.3333333333333335)--  (2,2);
					\draw [dashed]  (1,2.3333333333333335)-- (0,0);
					 \draw [fill] (0,0) circle (1pt);
					\draw[color=black] (-0.39,0.17)  node {$A$};
					\draw [fill] (-1,-1) circle  (1pt);
					 \draw[color=black] (-1.29,-1.13) node {$B$};
					\draw  [fill] (3,-1) circle (1pt);
					\draw[color=black]  (3.29,-1.01) node {$C$};
					\draw [fill]  (4,0) circle (1pt);
					 \draw[color=black] (4.17,0.29) node {$D$};
					\draw [fill] (0,4) circle  (1pt);
					\draw[color=black] (-0.03,4.45)  node {$S$};
					\draw [fill] (-0.5,1.5)  circle (1pt);
					\draw[color=black]  (-0.89,1.55) node {$M$};
					\draw [fill]  (2,2) circle (1pt);
					 \draw[color=black] (2.17,2.29) node {$P$};
					\draw [fill]  (1,2.3333333333333335) circle (1pt);
					\draw[color=black]  (0.51,2.59) node {$N$};
				\end{tikzpicture}
		\end{center}
		Gọi $M, N, P$ lần lượt là giao  điểm của $(P)$ với các đường thẳng $SB, SC, SD$. \\
		Ta có  $SA\perp (ABCD)\Rightarrow SA\perp BC$. Mà $BC\perp AB$. \\
		$\Rightarrow BC\perp (SAB)\Rightarrow BC\perp AM$. \\ 
		Mặt  khác $SC\perp (P)\Rightarrow SC\perp AM$ nên $AM\perp (SBC)\Rightarrow AM\perp SB$. \\
		Tương tự $AN\perp SC,\  AP\perp SD,\ MP\parallel BD\Rightarrow MP\perp AN$. \\
		Ta  có $\dfrac{SM}{SB}=\dfrac{SP}{SD}=\dfrac{MP}{BD}=\dfrac{4}{5}\Rightarrow MP=\dfrac{4a\sqrt{2}}{5}$. \\
		$\triangle SAN$ vuông tại $A$ nên $AN=\dfrac{AS\cdot  AC}{\sqrt{AS^2+AC^2}}=\dfrac{2a\sqrt{3}}{3}$. \\
		 Suy ra $S_{AMNP}=\dfrac{AN\cdot MP}{2}=\dfrac{4a^2\sqrt{6}}{15}$.
	 }
\end{ex}

% nguồn: Ok_TuanLuan_Toan11_Lan2
\begin{ex}%[1H8V4-4]%[ma:31a5b8ac]
	Cho lăng trụ tứ giác đều $ABCD.A'B'C'D'$ có cạnh đáy bằng $a$. Mặt phẳng $(\alpha)$ hợp với mặt phẳng đáy $(ABCD)$ một góc bằng $45^{\circ}$ và cắt các cạnh bên của lăng  trụ tại các điểm $M$, $N$, $P$, $Q$. Tính diện tích của thiết diện.
	 \loigiai{
	\immini{ 
		Gọi $S$ là diện  tích thiết diện $MNPQ$.\\
		Ta có hình chiếu của  $(MNPQ)$ xuống $(ABCD)$ chính là hình vuông $ABCD$.\\
		Ta có $S_{ABCD}=a^2$. Gọi $\theta=\left((\alpha); (ABCD)\right)$.\\ 
		 Khi đó $\theta=45^{\circ}$.\\
		Suy ra $S_{MNPQ}=S_{ABCD}\cdot\cos{\theta}=a^2\cdot\cos{45^{\circ}}=a^2\sqrt{2}$. 
		}{ 
		\begin{tikzpicture}[scale=1,font=\footnotesize,line join=round,line cap=round,>=stealth]
				\def\bc{5} 
				\def\ba{2} 
				\def\h{3} 
				 \def\gocB{30}
				 \coordinate (B) at (0,0);
				 \coordinate (A) at (\gocB:\ba);
				\coordinate  (C) at (\bc,0);
				 \coordinate (D) at ($(C)-(B)+(A)$);
				\coordinate  (A') at ($(A)+(90:\h)$);
				\coordinate (B') at ($(B)-(A)+(A')$); 
				\coordinate  (C') at ($(C)-(A)+(A')$);
				\coordinate (D')  at ($(D)-(A)+(A')$);
				\coordinate (M)  at ($(A)!1/3!(A')$);
				\coordinate  (N) at ($(B)!1/2!(B')$);
				\coordinate (P) at ($(C)!1/2!(C')$); 
				\coordinate (Q)  at ($(D)!1/3!(D')$);
				\draw  (B')--(B)--(C)--(D)--(D')--(A')--(B')--(C')--(D')
					 (C)--(C')
					 (N)--(P)--(Q);
				\draw[dashed]  (A')--(A)--(D)
					(A)--(B) 
					(N)--(M)--(Q); 
				\foreach \x/\g  in {A/-90,B/-90,C/-90,D/-90,A'/90,B'/90,C'/90,D'/90,M/135,N/135,P/135,Q/135} \fill (\x) circle (1pt)+(\g:3mm) node {$\x$};
			\end{tikzpicture}
		 }
	} 
\end{ex}

% nguồn: Ok_TuanLuan_Toan11_Lan2
\begin{ex}%[1H8V4-4]%[ma:c51b6415]
	Cho hình chóp $S.ABC$ có đáy $ABC$ là tam giác đều cạnh $a$. Họi $H$ là trung điểm của cạnh $BC$, $O$ là trung điểm của $AH$ và $G$ là trọng tâm của tam giác $ABC$. Biết $SO$ vuông  góc với mặt phẳng $(ABC)$ và $SO=2a$. Tính diện tích của thiết diện tạo bởi mặt phẳng $(P)$ đi qua $G$ và vuông góc với $AH$ với hình chóp.
	\loigiai{ 
		\begin{center}
			\begin{tikzpicture}[scale=1,font=\footnotesize,line join=round,line cap=round,>=stealth]
					\coordinate  (A) at (0,0);
					\coordinate (B) at  (5.5,-3);
					\coordinate  (C) at (8,0);
					\coordinate (H)  at ($(B)!1/2!(C)$);
					 \coordinate (G) at ($(A)!2/3!(H)$);
					\coordinate  (O) at ($(A)!1/2!(H)$);
					\coordinate (S) at  ($(O)+(0,6)$);
					 \coordinate (M) at ($(A)!2/3!(B)$);
					\coordinate (N)  at ($(A)!2/3!(C)$);
					\coordinate (K) at  ($(H)!2/3!(S)$);
					\coordinate (Q) at ($(B)!2/3!(S)$); 
					\coordinate (P) at ($(C)!2/3!(S)$); 
					\draw  (S)--(A)--(B)--(C)--(S)--(B)
						(S)--(H) 
						 (M)--(Q)--(P);
					\draw [dashed] (S)--(O) 
						 (C)--(A)--(H)
						 (M)--(N)--(P)
						(G)--(K); 
					\foreach \x/\g in {S/90,A/180,B/-90,C/0,O/-90,G/-90,M/-90,N/-90,H/-90,K/75,P/90,Q/180} \fill  (\x) circle (1pt)+(\g:3mm) node {$\x$};
				\end{tikzpicture}
		\end{center}
		Qua $G$ dựng đường thẳng $MN$  song song với $BC$, dựng $GK$ song song với $SO$.\\
		Khi  đó $MN\perp AH$, $GK\perp AH$.\\
		Suy  ra $MN\subset (P)$, $GK\subset (P)$.\\
		Qua điểm $K$ dựng đường thẳng $PQ$ song song với $BC$.\\ 
		 Khi đó $PQ\perp AH\Rightarrow PQ\subset (P)$.\\
		Thiết  diện là tứ giác $MNPQ$.\\
		Do $MN$ và $PQ$ cùng song song với $BC$ nên suy  ra $G$ là trung điểm của $MN$ và $K$ là trung điểm của $PQ$.\\
		Tứ giác $MNPQ$ là hình  thang.\\
		Vì $G$ là trọng tâm $\triangle{ABC}$ và $MN\parallel BC$  nên suy ra $MN=\dfrac{2}{3}BC=\dfrac{2a}{3}$.\\
		Ta có $\dfrac{OH}{AH}=\dfrac{1}{2}$, $\dfrac{HG}{AH}=\dfrac{1}{3}\Rightarrow\dfrac{HG}{OH}=\dfrac{2}{3}$.\\ 
		Vì  $GK\parallel SO$ nên $\dfrac{HG}{HO}=\dfrac{GK}{SO}=\dfrac{HK}{HS}=\dfrac{2}{3}\Rightarrow KG=\dfrac{4}{3}a$.\\
		Mặt khác $PQ\parallel  BC$, $\dfrac{HK}{HS}=\dfrac{2}{3}\Rightarrow\dfrac{SK}{SH}=\dfrac{1}{3}=\dfrac{PQ}{BC}\Rightarrow PQ=\dfrac{1}{3}a$.\\
		Vậy  diện tích thiết diện cần tìm là \[S=\dfrac{1}{2}(PQ+MN)\cdot GK=\dfrac{1}{2}\left(\dfrac{1}{3}a+\dfrac{2}{3}a\right)\cdot\dfrac{4}{3}a=\dfrac{2a^3}{3}.\]
	 }
\end{ex}
